Tarkan kohteentunnistusmallin opettaminen vaatii suuria määriä erilaisia kuvia tunnistettavasta kohteesta. Välttämättä aina ei ole saatavilla valmiita kuvia mallin opettamista varten, jos kyseessä on esimerkiksi perinteisistä värikuvista poikkeava esitysmuoto tai jos tunnistettava kohde ei ole yleinen. Tällöin mallin opettamiseen on otettava kuvat itse, mikä on työlästä ja aikaavievää. Perinteisesti tämän ongelman helpottamiseksi opetukseen käytettävien kuvien määrää on voitu kasvattaa augmentoimalla jo otettuja kuvia.

Nykyiset realistisia synteettisiä kuvia luovat tekstistä-kuvaksi -mallit ovat kehittyneet vauhdilla viimeisten vuosien aikana. Avoimella lähdekoodillaan suosituksi tullut Stable Diffusionin tekstistä-kuvaksi -malli SDXL mahdollistaa omien hienosäädettyjen synteettisiä kuvia tuottavien mallien luomisen. Jo pienellä määrällä kuvia voidaan kuvia generoiva malli hienosäätää tuottamaan kuvasetin sisältöä vastaavia synteettisiä kuvia. Parhaimmillaan nämä kuvat ovat tarpeeksi realistisia, jotta niillä voidaan kouluttaa kohteentunnistusmalli käytettäväksi reaalimaailman kuville. Tässä diplomityössä tutkitaan, voidaanko tuotetuilla synteettisillä kuvilla parantaa kohteentunnistusmallin suorituskykyä.

Työssä tuotettiin kohteentunnistusmalli ihmisten paikantamiseen lähi-infrapunakuvista pienissä sisätiloissa. Mallin kouluttamista varten tuotettiin lähi-infrapunakuvien kaltaisia synteettisiä kuvia SDXL tekstistä-kuvaksi -mallilla hienosäätämällä sitä 31 lähi-infrapunakuvan ja LoRA-tekniikan avulla. Näitä synteettisiä kuvia käytettiin yhdessä yhdessä aitojen kuvien kanssa luomaan joukko hienosäädettyjä YOLOv8n-kohteentunnistusmalleja, joiden suorituskykyä verrattiin keskenään tuloksissa.

Tuloksista voitiin havaita, että jo pelkästään synteettisillä kuvilla hienosäädetyt mallit suoriutuivat paremmin (AP\textsubscript{IoU:0.50}: 0.866-0.908) kuin perusmalli (AP\textsubscript{IoU:0.50}: 0.815). Kun reaalimaailman kuvien osuus kasvatettiin noin kymmenesosaan suuresta määrästä synteettisiä kuvia, kasvoi mallin suorituskyky ja stabiilisuus entisestään. Paras tulos saatiin, kun kaikki 568 reaalimaailman kuvaa yhdistettiin 2500 synteettiseen kuvaan (AP\textsubscript{IoU:0.50}: 0.938). Tämä työ osoitti, että hienosäädetyn diffuusiopohjaisen tekstistä-kuvaksi -mallin tuottamilla synteettisillä kuvilla voidaan parantaa kohteentunnistusmallin suorituskykyä tunnistamaan kohteita reaalimaailman kuvista.